% !TeX program = xelatex
\documentclass[11pt]{article}
\newcommand{\notesdir}{../../lecture-notes}
\providecommand{\notesdir}{.}

% -----------------------------------------------------------------------------
% Shared preamble for the CO5059 lecture notes and the quiz-review documents.
% Documents outside this folder define \notesdir (path to this folder) first.
% -----------------------------------------------------------------------------

% Page layout and typography
\usepackage[a4paper,portrait,margin=1in]{geometry}
\usepackage{fontspec}
\newfontfamily\arialfont{Arial}
\usepackage{microtype}
\usepackage{parskip}

% Mathematics and diagrams
\usepackage{amsmath,amssymb,amsthm,mathtools}
\usepackage{tikz-cd}
\usetikzlibrary{arrows.meta,positioning,shapes.geometric,shapes.arrows,calc,fit,backgrounds}

% Reusable styles for flowcharts and concept diagrams.
\tikzset{
  concept/.style={
    draw=brandprimary,
    fill=brandsecondary!8,
    rounded corners,
    minimum width=3.2cm,
    minimum height=0.9cm,
    align=center,
    font=\small,
  },
  diagram arrow/.style={-{Latex[length=2.5mm]},thick,brandprimary},
}

% Tables, figures, code, and lists
\usepackage{graphicx}
\graphicspath{{\notesdir/}}
\usepackage{enumitem}
\usepackage{booktabs}
\usepackage{array}
\usepackage{listings}

% Color and links
\usepackage[dvipsnames]{xcolor}
\definecolor{brandprimary}{HTML}{030391}
\definecolor{brandsecondary}{HTML}{1488D8}
\usepackage[
  colorlinks=true,
  linkcolor=brandprimary,
  citecolor=brandsecondary,
  urlcolor=brandsecondary
]{hyperref}

% Headers and footers
\usepackage{fancyhdr}
\setlength{\headheight}{36pt}
\pagestyle{fancy}
\fancyhf{}
\fancyhead[L]{%
  \parbox[c][0.42in][c]{0.42\textwidth}{%
    \raggedright\small\textbf{\coursecode}\\\coursename}}
\fancyhead[R]{%
  \parbox[c][0.42in][c]{0.42\textwidth}{%
    \raggedleft\includegraphics[height=0.42in]{assets/hcmut_logo.png}}}
\fancyfoot[L]{\small Semester \semester}
\fancyfoot[C]{\small\nouppercase{\leftmark}}
\fancyfoot[R]{\small Page \thepage}
\renewcommand{\sectionmark}[1]{\markboth{#1}{}}

% -----------------------------------------------------------------------------
% Course information -- change these values once per course
% -----------------------------------------------------------------------------
\newcommand{\universityname}{HCMC University of Technology}
\newcommand{\faculty}{Faculty of Computer Science and Engineering}
\newcommand{\coursecode}{CO5059}
\newcommand{\coursename}{Advanced System Architectures}
\newcommand{\semester}{261}
\newcommand{\instructor}{Tran Ngoc Thinh}
\newcommand{\instructoremail}{}% leave empty to hide
\newcommand{\notetaker}{Phan Nhat Minh}
\newcommand{\notetakeremail}{minh.phannhat95@hcmut.edu.vn}
\newcommand{\coursedescription}{Quantitative design and evaluation of computer
systems: performance and cost, instruction set architectures, pipelining and
its hazards, memory hierarchy, multiprocessors, superscalar/VLIW processors,
and AI \& security chips.}
\newcommand{\emaillink}[1]{%
  \href{mailto:#1}{\textcolor{brandprimary}{\texttt{#1}}}}

% Colored definitions, results, examples, proofs, and notes.
\input{\notesdir/colored-boxes}

% -----------------------------------------------------------------------------
% Course-specific notation
% -----------------------------------------------------------------------------
\newcommand{\CPI}{\mathrm{CPI}}

\lstdefinestyle{mips}{
  basicstyle=\ttfamily\small,
  keywordstyle=\bfseries\color{brandprimary},
  morekeywords={lw,sw,add,sub,and,or,slt,beq,bne,beqz,xor,nop,addi,addiu,addu,
    subu,li,la,jal,jr,jalr,sll,srl,sra,sllv,srlv,srav,andi,ori,xori,nor,slti,
    sltiu,sltu,lui,mult,multu,div,divu,mfhi,mflo,ll,sc,syscall,move,lh,lhu,lb,
    lbu,sh,sb,blez,bgtz,bltz,bgez,j,mul},
  showstringspaces=false,
  columns=fullflexible,
  keepspaces=true,
}
\lstset{style=mips}

% Creates a consistently formatted lecture heading.
% Usage: \lecture{number}{title}{date}
\newcommand{\lecture}[3]{%
  \section{Lecture #1: #2}
  \label{lec:#1}
  \begin{center}
    \small\textit{#3}
  \end{center}
}


% Cover page + contents page.  Usage: \makecover{Big title}{Description}
\newcommand{\makecover}[2]{%
\pagenumbering{roman}
\begin{titlepage}
  \thispagestyle{empty}

  \begin{center}
    \includegraphics[width=0.82\textwidth]{assets/hcmut_banner_en.png}

    \vspace{0.35cm}
    {\arialfont\large\bfseries\color{brandprimary}\faculty\par}
  \end{center}

  \vfill
  \begin{center}
    {\Huge\bfseries #1\par}
    \vspace{0.55cm}
    {\LARGE\bfseries\color{brandprimary}\coursename\par}
    \vspace{0.2cm}
    {\arialfont\large\bfseries\color{brandsecondary}\coursecode\par}

    \vspace{0.65cm}
    {\color{brandsecondary}\rule{0.38\textwidth}{1.2pt}}
    \vspace{0.65cm}

    {\large\textbf{Instructor:} \instructor
      \ifx\instructoremail\empty\else\ {\small(\emaillink{\instructoremail})}\fi\par}
    \vspace{0.3cm}
    {\large\textbf{Notes by:} \notetaker\
      {\small(\emaillink{\notetakeremail})}\par}
    \vspace{0.3cm}
    {\large\textbf{Semester:} \semester\par}

    \vspace{0.8cm}
    \begin{minipage}{0.72\textwidth}
      \centering\itshape\color{black!70}
      #2
    \end{minipage}
  \end{center}

  \vfill
  \begin{center}
    \small\color{black!60}Last updated: \today
  \end{center}
\end{titlepage}

% -----------------------------------------------------------------------------
% Table of contents -- kept on its own page
% -----------------------------------------------------------------------------
\clearpage
\setcounter{page}{2}
\thispagestyle{empty}
\hypersetup{linkcolor=black}
\setcounter{tocdepth}{2}
\tableofcontents
\hypersetup{linkcolor=brandprimary}
\clearpage
\pagenumbering{arabic}
}


\begin{document}
\sloppy
\makecover{Quiz Review: Pipelining}{Lecture 3 -- pipeline timing, speedup,
CPI with stalls, structural/data/control hazards, forwarding, branch
schemes. A cheat sheet, the slide exercises with solutions, exam-style
questions, and extra practice problems -- with solutions at the end.}

% ============================================================================
\section{Cheat Sheet}

\subsection{Timing}
\begin{theorem}[Cycles and time]
For $n$ instructions on a $k$-stage pipeline with $s$ total stall cycles:
\[
  \text{Cycles} = k + (n-1) + s, \qquad
  \text{Time} = \text{Cycles}\times t_{clk},
\]
\[
  t_{clk} = \max_i(\text{stage delay}_i) + \text{latch overhead},
  \qquad
  \text{Latency of one instruction} = k\cdot t_{clk}.
\]
\end{theorem}

\begin{theorem}[Speedup]
Ideal (balanced stages, $p$ = unpipelined time per task):
\[
  \text{Speedup} = \frac{n\,p}{\frac{p}{k}(n-1)+p} = \frac{nk}{n+k-1}
  \;\xrightarrow{n\to\infty}\; k.
\]
With stalls:
\begin{gather*}
  \CPI_{pipe} = 1 + \text{stall cycles per instruction}\\
  \text{Speedup} = \frac{\CPI_{unpipe}}{\CPI_{pipe}}\times
  \frac{t_{clk,unpipe}}{t_{clk,pipe}}
  = \frac{\text{Depth}}{1+\text{stall CPI}} \ \ (\text{balanced stages})
\end{gather*}
\end{theorem}

\begin{note}
Pipelining improves \textbf{throughput}, not \textbf{latency}. Ideal
$\CPI=1$ because one instruction \emph{completes} per cycle once the
pipeline is full: $\CPI = (k+n-1)/n \to 1$.
\end{note}

\subsection{The 5-stage MIPS pipeline}
IF (instruction memory) -- ID (decode, read registers) -- EX (ALU / address)
-- MEM (data memory) -- WB (write register). Registers are
\textbf{written in the first half} and \textbf{read in the second half} of a
cycle, so an instruction can read in ID a value written by WB in the same
cycle.

\subsection{Data hazards: how many stalls?}
Let $d$ be the distance from producer to consumer ($d=1$: the very next
instruction).
\begin{center}
\begin{tabular}{@{}p{4.6cm}lp{6.2cm}@{}}
\toprule
Situation & Stalls & Why \\
\midrule
No forwarding, any producer & $\max(0,\,3-d)$ & consumer's ID must not precede producer's WB \\
Forwarding, ALU $\to$ ALU & 0 & EX/MEM $\to$ EX ($d=1$), MEM/WB $\to$ EX ($d=2$) \\
Forwarding, load $\to$ use ($d=1$) & 1 & data ready only at end of MEM \\
Forwarding, load $\to$ use ($d\ge2$) & 0 & MEM/WB $\to$ EX \\
Forwarding, load $\to$ store data & 0 & MEM/WB $\to$ store input \\
\bottomrule
\end{tabular}
\end{center}
Only \textbf{RAW} causes hazards in the in-order 5-stage pipeline; WAR and
WAW cannot happen because reads (ID) are early and writes (WB) are late and
in order.

\subsection{Drawing a pipeline diagram}
\begin{enumerate}[leftmargin=*]
  \item One row per instruction, one column per cycle; each row starts one
  cycle after the previous row's IF.
  \item When an instruction stalls, write \texttt{stall} after its ID (it
  waits in ID); the instruction behind it waits in IF, and everything behind
  shifts right.
  \item Total cycles $=$ cycle of the last WB.
\end{enumerate}
In wide diagrams below, \texttt{stl} abbreviates \texttt{stall}.

\subsection{Control hazards}
\begin{prop}[Branch cost]
\begin{gather*}
  \text{Branch penalty} = (\text{stage where branch is resolved}) - 1\\
  \CPI = 1 + f_{branch}\times(\text{fraction penalized})\times\text{penalty}
\end{gather*}
Resolved in MEM $\Rightarrow$ 3 cycles; moved to ID (zero test + extra adder)
$\Rightarrow$ 1 cycle.
\end{prop}
\begin{center}
\small
\begin{tabular}{@{}lp{9.6cm}@{}}
\toprule
Scheme & Cost \\
\midrule
Stall (flush) & every branch pays the full penalty \\
Predict not taken & only \emph{taken} branches pay (fetch PC$+4$, squash on taken); $\approx47\%$ of MIPS branches are not taken \\
Predict taken & untaken branches pay; in the 5-stage MIPS still 1 cycle because the target is not known earlier \\
Delayed branch & the instruction in the delay slot always executes; cost $=$ wasted (unfilled/useless) slots \\
\bottomrule
\end{tabular}
\end{center}
Delay-slot filling: (A) from before the branch -- best, always useful, the
branch must not depend on it; (B) from the target -- useful when taken, may
need copying, must be safe if not taken; (C) from the fall-through -- useful
when not taken, must be safe if taken. Compilers fill $\approx60\%$ of slots,
$\approx80\%$ of those usefully $\Rightarrow\approx50\%$ useful.

\subsection{Structural hazards}
Two instructions need the same resource in the same cycle (e.g.\ one memory
port for instructions and data). Fix: duplicate the resource (split I/D
memory) or stall. Sometimes accepted when the hazard is rare and avoiding it
is costly.

% ============================================================================
\section{Exercises from the Slides (with Solutions)}

\begin{exercise}[Laundry]
Wash 30~min, dry 40~min, fold 20~min; 4 loads. Sequential vs.\ pipelined
time and speedup?
\end{exercise}
\begin{solution}
Sequential: $4\times90 = 360$~min $=6$~h. Pipelined: the 40-min dryer sets
the pace: $30 + 4\times40 + 20 = 210$~min $=3.5$~h. Speedup $=6/3.5 =
1.7$ (not 3), because stages are unbalanced and the pipeline must fill and
drain.
\end{solution}

\begin{exercise}[Single-cycle vs.\ pipelined]
Stage times IF 2, Reg 1, ALU 2, Mem 2, WB 1~ns. Compare 3 loads, then 1003
loads, on a single-cycle (8~ns) and a pipelined (2~ns) machine.
\end{exercise}
\begin{solution}
3 loads: single-cycle $3\times8 = 24$~ns; pipelined $(5+2)\times2 = 14$~ns;
speedup $1.7$.
1003 loads: single-cycle $1003\times8 = 8024$~ns; pipelined
$(5+1002)\times2 = 2014$~ns; speedup $3.98 \approx 8/2 = 4$. The speedup
approaches the ratio of clock periods (not 5, because the stages are
unbalanced: the 1-ns stages still take a 2-ns cycle).
\end{solution}

\begin{exercise}[Pipelining a multi-cycle CPU]
Unpipelined: 1~ns clock, ALU/branch 4 cycles, memory 5 cycles; mix 40\%
ALU, 20\% branch, 40\% memory. Pipelining adds 0.2~ns to the clock.
Speedup?
\end{exercise}
\begin{solution}
Unpipelined average $=1\times(0.6\times4 + 0.4\times5) = 4.4$~ns.
Pipelined average $= 1.2$~ns (ideal $\CPI=1$). Speedup $=4.4/1.2 = 3.7$.
\end{solution}

\begin{exercise}[Structural hazard]
Loads/stores are 40\% of instructions; the machine with one memory port
stalls 1 cycle per data access but has a $1.05\times$ faster clock. Which is
better?
\end{exercise}
\begin{solution}
Hazard machine: $(1+0.4\times1)\times t/1.05 = 1.33\,t$ versus $1.0\,t$ for
the ideal machine, so the machine without the hazard is $1.33\times$
faster.
\end{solution}

\begin{exercise}[RAW hazard on R1]
\texttt{add r1,r2,r3; sub r4,r1,r3; and r6,r1,r7; or r8,r1,r9;
xor r10,r1,r11}. Stalls and total cycles without and with forwarding?
\end{exercise}
\begin{solution}
Without forwarding \texttt{sub} ($d=1$) needs 2 stalls; after that every
later reader is far enough behind:
@@GRID:raw_nofwd@@
11 cycles. With forwarding (EX/MEM $\to$ EX to \texttt{sub}, MEM/WB $\to$ EX
to \texttt{and}, register file for \texttt{or}/\texttt{xor}): 0 stalls,
$5+4 = 9$ cycles.
\end{solution}

\begin{exercise}[Load-use hazard]
\texttt{lw r1,0(r2); sub r4,r1,r6; and r6,r1,r7; or r8,r1,r9} with
forwarding.
\end{exercise}
\begin{solution}
The loaded value exists only at the end of MEM (cycle 4) but \texttt{sub}
needs it at the start of its EX (cycle 4): 1 stall is unavoidable.
@@GRID:loaduse@@
9 cycles ($5+3+1$).
\end{solution}

\begin{exercise}[Compiler scheduling]
\texttt{a=b+c; d=e-f;} compiled as \texttt{lw Rb; lw Rc; add Ra; sw Ra;
lw Re; lw Rf; sub Rd; sw Rd} (forwarding, 1-cycle load delay). Stalls
before and after scheduling?
\end{exercise}
\begin{solution}
Original: \texttt{add} right after \texttt{lw Rc} and \texttt{sub} right
after \texttt{lw Rf} $\Rightarrow$ 2 stalls, $5+7+2 = 14$ cycles.
Scheduled \texttt{lw Rb; lw Rc; lw Re; add Ra; lw Rf; sw Ra; sub Rd; sw Rd}:
no load is immediately followed by its use $\Rightarrow$ 0 stalls, 12
cycles.
\end{solution}

\begin{exercise}[Branch stall impact]
30\% of instructions are branches with a 3-cycle stall. New CPI?
\end{exercise}
\begin{solution}
$\CPI = 1 + 0.30\times3 = 1.9$ -- almost twice the ideal. With the branch
resolved in ID (1 cycle): $1+0.3 = 1.3$.
\end{solution}

\begin{exercise}[Comparing branch schemes]
4\% unconditional, 6\% conditional untaken, 10\% conditional taken;
pipeline depth 5. Find CPI and speedups for stall, predict not taken,
predict taken, delayed branch (penalty 0.5).
\end{exercise}
\begin{solution}
\begin{tabular}[t]{@{}lcccc@{}}
Scheme & CPI & vs.\ unpipelined & vs.\ stall \\
\hline
Stall: $1+0.20\times3$ & 1.60 & $5/1.6=3.1$ & 1.00 \\
Not taken: $1+0.04\times1+0.10\times3$ & 1.34 & 3.7 & 1.19 \\
Taken: $1+0.14\times1+0.06\times2$ & 1.26 & 4.0 & 1.27 \\
Delayed: $1+0.20\times0.5$ & 1.10 & 4.5 & 1.45 \\
\end{tabular}

(The slide lists 1.29 for predict taken, which is the ratio of the rounded
speedups $4.0/3.1$; the exact value is $1.6/1.26 = 1.27$.)
\end{solution}

\begin{exercise}[CPI of a realistic mix]
Loads 30\% (25\% followed immediately by a use), branches 20\% (45\%
taken); forwarding, branch resolved in ID, predict not taken.
\end{exercise}
\begin{solution}
$\CPI = 1 + 0.3\times0.25\times1 + 0.2\times0.45\times1 = 1.165$.
\end{solution}

\begin{exercise}[Hazards, NOPs and forwarding]\label{e:lecture}
\leavevmode
\begin{lstlisting}
(1) lw  $1, 40($6)
(2) add $6, $2, $2
(3) sw  $6, 50($1)
(4) add $4, $5, $6
(5) lw  $6, 10($4)
\end{lstlisting}
(A) Find all data dependencies. (B) Without forwarding, which cause hazards?
(C) Insert NOPs so the code runs correctly without forwarding. (D) How many
cycles does (C) take? (E) With full forwarding, how many cycles, and which
forwarding paths are used?
\end{exercise}
\begin{solution}
\textbf{(A)} RAW: (1)$\to$(3) on \texttt{\$1}; (2)$\to$(3) on \texttt{\$6};
(2)$\to$(4) on \texttt{\$6}; (4)$\to$(5) on \texttt{\$4}.
WAR: (1)$\to$(2), (3)$\to$(5), (4)$\to$(5) on \texttt{\$6}.
WAW: (2)$\to$(5) on \texttt{\$6}.

\textbf{(B)} Only RAW with $d\le2$: (1)$\to$(3) ($d=2$), (2)$\to$(3)
($d=1$), (2)$\to$(4) ($d=2$), (4)$\to$(5) ($d=1$). WAR/WAW are not hazards
here.

\textbf{(C)} Each consumer must be at least 3 slots after its producer:
\begin{lstlisting}
lw  $1, 40($6)
add $6, $2, $2
nop
nop
sw  $6, 50($1)
add $4, $5, $6
nop
nop
lw  $6, 10($4)
\end{lstlisting}
\textbf{(D)} 9 instructions $\Rightarrow 5 + 8 = 13$ cycles:
@@GRID:lec_nofwd@@

\textbf{(E)} No stalls are needed (the only load, (1), is used two
instructions later), so $5+4 = 9$ cycles:
@@GRID:lec_fwd@@
Forwarding paths: cycle 5, MEM/WB $\to$ EX gives \texttt{\$1} (from (1)) and
EX/MEM $\to$ EX gives \texttt{\$6} (from (2)) to the \texttt{sw}; cycle 6,
MEM/WB $\to$ EX gives \texttt{\$6} (from (2)) to (4); cycle 7, EX/MEM $\to$
EX gives \texttt{\$4} (from (4)) to (5).
\end{solution}

% ============================================================================
\section{Exam-Style Questions}

Multi-part questions in the style of the in-class exercise. Solutions are in
the last section.

\begin{exercise}[Pipeline timing]\label{x:timing}
A 5-stage pipeline has stage latencies IF 250, ID 200, EX 300, MEM 350,
WB 150~ps. Each pipeline register adds 20~ps.
\begin{enumerate}[label=(\alph*),leftmargin=*]
  \item What is the clock period of the non-pipelined (single-cycle) version?
  \item What is the clock period of the pipelined version?
  \item What is the latency of one instruction in the pipeline?
  \item How long do $10^6$ instructions take on each version (no stalls)?
  What is the speedup?
  \item What is the pipelined throughput in instructions per second?
  \item MEM is split into two stages of 175~ps each. Find the new clock period
  and the speedup over the 5-stage pipeline for $10^6$ instructions.
\end{enumerate}
\end{exercise}

\begin{exercise}[Hazards and forwarding]\label{x:hazards}
\leavevmode
\begin{lstlisting}
(1) lw  $s2, 0($s1)
(2) add $s3, $s2, $s4
(3) sub $s5, $s3, $s2
(4) sw  $s5, 4($s1)
(5) add $t0, $s6, $s7
\end{lstlisting}
\begin{enumerate}[label=(\alph*),leftmargin=*]
  \item List all RAW dependencies, with their distance.
  \item Which of them cause hazards without forwarding?
  \item Insert NOPs so the code runs correctly without forwarding. How many
  cycles does it take?
  \item With full forwarding, draw the pipeline diagram. How many stalls and
  cycles?
  \item Name the forwarding path used for every dependency in (d).
  \item Reorder the code (same result) so it runs with no stalls when
  forwarding is available. How many cycles now?
\end{enumerate}
\end{exercise}

\begin{exercise}[Branches and CPI]\label{x:branches}
A program's mix: 25\% loads, 10\% stores, 15\% branches (60\% taken), 50\%
ALU. 40\% of loads are immediately followed by an instruction that uses the
loaded value. The pipeline has full forwarding.
\begin{enumerate}[label=(\alph*),leftmargin=*]
  \item Branches are resolved in MEM and the pipeline stalls on every branch.
  CPI?
  \item Predict not taken, branch resolved in MEM. CPI?
  \item Predict not taken, branch resolved in ID. CPI?
  \item One delay slot (branch resolved in ID); the compiler fills 70\% of
  slots with useful instructions. CPI?
  \item With a 1~GHz clock and $10^9$ instructions, find the execution time
  for (a) and (c) and the speedup.
  \item An unpipelined multi-cycle version runs at the same clock with
  $\CPI = 4.6$. What speedup does design (c) give over it?
\end{enumerate}
\end{exercise}

% ============================================================================
\section{Practice Problems}

\begin{exercise}\label{p:clock}
Stage delays: IF 200, ID 150, EX 250, MEM 300, WB 150~ps; latch overhead
20~ps. (a) Single-cycle clock? (b) Pipelined clock? (c) Latency of one
instruction when pipelined? (d) Time for 1000 instructions on each, and the
speedup? (e) MEM is split into two 150-ps stages. New clock and speedup over
(d) for 1000 instructions?
\end{exercise}

\begin{exercise}\label{p:ideal}
For a 5-stage pipeline, compute the ideal speedup for $n=10$, $n=100$ and
$n\to\infty$ instructions.
\end{exercise}

\begin{exercise}\label{p:stalls}
\leavevmode
\begin{lstlisting}
(1) add $t0, $t1, $t2
(2) lw  $t3, 0($t0)
(3) sub $t4, $t3, $t5
(4) sw  $t4, 4($t0)
(5) or  $t6, $t0, $t4
\end{lstlisting}
List the RAW dependencies. Draw the diagram and count cycles (a) without
forwarding, (b) with full forwarding.
\end{exercise}

\begin{exercise}\label{p:sched}
With forwarding, how many cycles does this take? Reorder it (without changing
the result) to remove all stalls.
\begin{lstlisting}
lw  $t0, 0($a0)
add $t1, $t0, $t2
lw  $t3, 4($a0)
add $t4, $t3, $t2
sw  $t1, 0($a1)
\end{lstlisting}
\end{exercise}

\begin{exercise}\label{p:branch}
25\% of instructions are branches and 60\% of them are taken. Find the CPI
for: (a) stall, branch resolved in MEM; (b) predict not taken, resolved in
MEM; (c) predict not taken, resolved in ID; (d) one delay slot that is
usefully filled 50\% of the time.
\end{exercise}

\begin{exercise}\label{p:mix}
Loads 30\% (40\% of them followed immediately by a use), branches 15\% (50\%
taken, 1-cycle penalty on taken), full forwarding. Find the CPI. An
unpipelined version has CPI 4.5 with a 1~ns clock; the pipelined clock is
1.1~ns. Speedup?
\end{exercise}

\begin{exercise}\label{p:struct}
35\% of instructions access data memory. With a single memory port each
access causes a 1-cycle stall. A dual-port design avoids this but its clock
period is 10\% longer. Which is faster and by how much?
\end{exercise}

\begin{exercise}\label{p:classify}
Classify every dependency (RAW/WAR/WAW). Which ones are hazards in the
5-stage MIPS pipeline?
\begin{lstlisting}
(1) add $1, $2, $3
(2) sub $4, $1, $5
(3) and $1, $6, $7
(4) or  $8, $1, $4
\end{lstlisting}
\end{exercise}

\begin{exercise}\label{p:slot}
Can \texttt{add \$s1,\$s2,\$s3} be moved into the delay slot of the
following \texttt{beq \$t0,\$t1,L}? What if it were
\texttt{add \$t0,\$s2,\$s3} instead?
\end{exercise}

\begin{exercise}\label{p:draw}
Draw the diagram with full forwarding, mark the stall(s), and name the
forwarding path for every dependency:
\begin{lstlisting}
(1) lw  $2, 20($1)
(2) and $4, $2, $5
(3) or  $8, $2, $6
(4) add $9, $4, $2
(5) slt $1, $6, $7
\end{lstlisting}
\end{exercise}

% ============================================================================
\section{Solutions to the Practice Problems}

\begin{solution}[\ref{p:clock}]
(a) $200+150+250+300+150 = 1050$~ps.
(b) $300+20 = 320$~ps (slowest stage plus latch).
(c) $5\times320 = 1600$~ps -- longer than single-cycle; pipelining does not
cut latency.
(d) Single-cycle $1000\times1050 = 1{,}050{,}000$~ps; pipelined
$(5+999)\times320 = 321{,}280$~ps; speedup $3.27$.
(e) Slowest stage is now EX: $250+20=270$~ps; $(6+999)\times270 =
271{,}350$~ps; speedup over (d) $=321{,}280/271{,}350 = 1.18$.
\end{solution}

\begin{solution}[\ref{p:ideal}]
$nk/(n+k-1)$: $n=10$: $50/14 = 3.57$; $n=100$: $500/104 = 4.81$;
$n\to\infty$: 5.
\end{solution}

\begin{solution}[\ref{p:stalls}]
RAW: (1)$\to$(2), (1)$\to$(4), (1)$\to$(5) on \texttt{\$t0};
(2)$\to$(3) on \texttt{\$t3}; (3)$\to$(4), (3)$\to$(5) on \texttt{\$t4}.

(a) Without forwarding (each consumer $\ge3$ slots after its producer):
@@GRID:p3_nofwd@@
6 stalls, 15 cycles.

(b) With forwarding only the load-use (2)$\to$(3) stalls:
@@GRID:p3_fwd@@
1 stall, 10 cycles.
\end{solution}

\begin{solution}[\ref{p:sched}]
Two load-use stalls: $5+4+2 = 11$ cycles. Reordered:
\begin{lstlisting}
lw  $t0, 0($a0)
lw  $t3, 4($a0)
add $t1, $t0, $t2
add $t4, $t3, $t2
sw  $t1, 0($a1)
\end{lstlisting}
Every use is now 2 instructions after its load: 0 stalls, 9 cycles.
\end{solution}

\begin{solution}[\ref{p:branch}]
(a) $1+0.25\times3 = 1.75$.
(b) $1+0.25\times0.6\times3 = 1.45$.
(c) $1+0.25\times0.6\times1 = 1.15$.
(d) $1+0.25\times0.5 = 1.125$.
\end{solution}

\begin{solution}[\ref{p:mix}]
$\CPI = 1+0.3\times0.4\times1+0.15\times0.5\times1 = 1+0.12+0.075 = 1.195$.
Speedup $=(4.5\times1)/(1.195\times1.1) = 4.5/1.3145 = 3.42$.
\end{solution}

\begin{solution}[\ref{p:struct}]
Single port: $1.35\times t$. Dual port: $1.0\times1.1t = 1.1t$. Dual port is
faster by $1.35/1.1 = 1.23\times$.
\end{solution}

\begin{solution}[\ref{p:classify}]
RAW: (1)$\to$(2) \texttt{\$1}; (2)$\to$(4) \texttt{\$4}; (3)$\to$(4)
\texttt{\$1} (instruction 4 reads the \emph{new} \texttt{\$1} from (3), not
from (1)). WAR: (2)$\to$(3) on \texttt{\$1}. WAW: (1)$\to$(3) on
\texttt{\$1}. Only the RAW dependencies can be hazards in the in-order
5-stage pipeline; with forwarding none of them stalls.
\end{solution}

\begin{solution}[\ref{p:slot}]
Yes: the branch does not read \texttt{\$s1}, so it is an independent
instruction from before the branch (strategy A, always useful). No for
\texttt{add \$t0,\dots}: the branch depends on \texttt{\$t0}, so the slot
must be filled from the target (B) or fall-through (C) path instead.
\end{solution}

\begin{solution}[\ref{p:draw}]
@@GRID:p10@@
(1)$\to$(2) \texttt{\$2}: load-use, 1 stall, then MEM/WB $\to$ EX in
cycle 5. (1)$\to$(3) \texttt{\$2}: register file (written in the first half
of cycle 5, read in the second half). (2)$\to$(4) \texttt{\$4}: MEM/WB $\to$
EX in cycle 7. (1)$\to$(4) \texttt{\$2}: register file. (5) has no
dependency. Total 10 cycles.
\end{solution}

% ============================================================================
\section{Solutions to the Exam-Style Questions}

\begin{solution}[\ref{x:timing}]
\leavevmode
\begin{enumerate}[label=(\alph*),leftmargin=*]
  \item $250+200+300+350+150 = 1250$~ps.
  \item Slowest stage $+$ overhead $= 350+20 = 370$~ps.
  \item $5\times370 = 1850$~ps (longer than 1250~ps: latency gets worse).
  \item Non-pipelined: $10^6\times1250\text{ ps} = 1.25$~ms. Pipelined:
  $(5 + 10^6 - 1)\times370\text{ ps} \approx 0.370$~ms. Speedup $\approx
  1250/370 = 3.38$ (well below 5 because the stages are unbalanced).
  \item One instruction per cycle: $1/370\text{ ps} \approx 2.7\times10^9$
  instructions/s.
  \item The slowest stage is now EX: $300+20 = 320$~ps. Speedup $=
  \dfrac{(10^6+4)\times370}{(10^6+5)\times320} \approx 1.156$.
\end{enumerate}
\end{solution}

\begin{solution}[\ref{x:hazards}]
\leavevmode
\begin{enumerate}[label=(\alph*),leftmargin=*]
  \item (1)$\to$(2) \texttt{\$s2} ($d=1$); (1)$\to$(3) \texttt{\$s2} ($d=2$);
  (2)$\to$(3) \texttt{\$s3} ($d=1$); (3)$\to$(4) \texttt{\$s5} ($d=1$).
  Instruction (5) is independent.
  \item All four have $d\le2$, so all are hazards without forwarding.
  \item Two NOPs after each of (1), (2), (3):
  \texttt{lw; nop; nop; add; nop; nop; sub; nop; nop; sw; add} -- 11
  instructions, $5+10 = 15$ cycles. (With an interlock instead of NOPs the
  timing is the same:)
@@GRID:x2_nofwd@@
  \item Only the load-use (1)$\to$(2) stalls: 1 stall, 10 cycles.
@@GRID:x2_fwd@@
  \item (1)$\to$(2): MEM/WB $\to$ EX in cycle 5 (after the stall).
  (1)$\to$(3): register file (written in the first half of cycle 5, read in
  the second half). (2)$\to$(3): EX/MEM $\to$ EX in cycle 6. (3)$\to$(4):
  EX/MEM $\to$ EX in cycle 7 (store data).
  \item Move the independent (5) between the load and its use:
  \texttt{lw; add \$t0; add \$s3; sub; sw}. The use of \texttt{\$s2} is now
  2 instructions after the load, so there are no stalls: 9 cycles.
@@GRID:x2_sched@@
\end{enumerate}
\end{solution}

\begin{solution}[\ref{x:branches}]
Load-use stalls per instruction $= 0.25\times0.4\times1 = 0.10$.
\begin{enumerate}[label=(\alph*),leftmargin=*]
  \item $1 + 0.10 + 0.15\times3 = 1.55$.
  \item $1 + 0.10 + 0.15\times0.6\times3 = 1.37$.
  \item $1 + 0.10 + 0.15\times0.6\times1 = 1.19$.
  \item Wasted slots: $0.15\times0.3\times1 = 0.045$, so $\CPI = 1.145$.
  \item $T_a = 10^9\times1.55\times1\text{ ns} = 1.55$~s; $T_c = 1.19$~s;
  speedup $= 1.55/1.19 = 1.30$.
  \item $4.6/1.19 = 3.87$.
\end{enumerate}
\end{solution}

\end{document}
